\subsection{Instruction-Guided Video Editing}
\label{subsec:edit_background}
In the task of text-based video editing\footnote{For brevity, we refer to text or instruction-based video editing simply as video editing.} the user provides the model with a video (either real or generated) along with an editing instruction text that specifies how they would like to alter the video.
The model is then expected to precisely modify the input video according to the given instruction, changing only the specified elements while preserving those that should remain intact.
The main challenge in developing a high-performing video editing model arises from the difficulty in collecting supervised data for this task.

As a result of this challenge, most prior work relies on training-free approaches~\citep{meng2021sdedit,Geyer2023TokenFlowCD,vid2vid-zero,text2video-zero,li2023vidtome,Ceylan2023Pix2VideoVE,2312.04524,yang2023rerender}, which can be applied to any \textToV model without requiring additional training.
In contrast to training-free methods, some approaches train models to generate videos by providing additional features of the video as input (\eg, depth or segmentation maps)~\citep{esser2023structure,Liang2023FlowVidTI,Yan2023MotionConditionedIA}.
However, these methods are inherently limited, as they cannot control features that were not incorporated during training.
For example, preserving the identity of an object or subject is impossible when conditioning only on the depth maps of a video.
Overall, both training-free methods and feature-based approaches tend to be imprecise and have been shown to perform worse than methods that explicitly adapt model parameters to process the entire video input during training~\citep{EVE,Qin2023InstructVid2VidCV}.

The current state-of-the-art approach for video editing, Emu Video Edit (EVE)~\citep{EVE}, employs two training stages to develop a video editing model.
First, it trains dedicated adapters for text-image-to-video generation and image editing on top of a shared \textToI model.
Next, it performs Factorized Diffusion Distillation (FDD) to align the adapters towards video editing.
In each training step of FDD, the model first generates an edited video through multiple diffusion steps.
Then, the edited video is given as input to two adversarial losses and two knowledge distillation losses, which provide supervision for the quality of the edited video.
Finally, this supervision is backpropagated through both the different losses and the entire generation chain.

We identify several key differences when comparing our approach to EVE.
First, we initialize training from a \textToV model and perform full model training (\Cref{subsec:edit_stage_1}), rather than training adapters for \textToV and image editing on top of a shared \textToI model.
Additionally, EVE's FDD backpropagates supervision through multiple forward passes with the model during generation, and an additional forward pass using each of the models it uses to provide supervision.
This makes FDD an order of magnitude more memory demanding than our approach, which limits the scalability of FDD.
By applying our approach to the \OursVideo (see~\Cref{sec:image_video_model}, we demonstrate that we can significantly surpass the reported results by EVE and set new state-of-the-art results in video editing.
